\documentclass{article}
\usepackage[margin=1.5cm]{geometry}
\usepackage{amsmath}
\usepackage{enumitem}

\setlist{nosep,leftmargin=*}

\begin{document}
\twocolumn

\title{Chapter 9 In-Class Exercise: Course Schedules and Modules}
\author{Prof. Jordan C. Hanson}
\date{}

\maketitle

\section{Adding Meeting Times}

\begin{enumerate}

\item Extend the course records from the warm-up by adding the meeting days and meeting time:

\begin{verbatim}
[subject,number,building,room,days,time]
\end{verbatim}

For example,

\begin{verbatim}
course = ["PHYS",100,"SLC",202,"MWF",1400]
\end{verbatim}

represents Physics 100, meeting in SLC 202 on Monday, Wednesday, and Friday at 1400 (2:00 PM).  Military time is useful because numerical times can be sorted directly:

\begin{verbatim}
930 < 1100 < 1400 < 1530
\end{verbatim}

\end{enumerate}


\section{Build the Module}

Create a new Python file called \verb+course_tools.py+.  All of our course-management functions will be placed inside
this module.  Begin with the function from the warm-up:

\begin{verbatim}
def add_course(courses, course):
    courses.append(course)
\end{verbatim}

Now define a function that returns the meeting time:

\begin{verbatim}
def get_time(course):
    return course[5]
\end{verbatim}

Use it to sort the master list:

\begin{verbatim}
def sort_by_time(courses):
    courses.sort(key=get_time)
\end{verbatim}

The \texttt{key} tells Python which part of each record should control the sorting.

\section{Displaying the Schedule}

Add one final function:

\begin{verbatim}
def display_courses(courses):
    for course in courses:
        print(course[0],
              course[1],
              course[4],
              f"{course[5]:04d}",
              course[2],
              course[3])
\end{verbatim}

The expression \verb+f"{course[5]:04d}"+ ensures that a time such as \texttt{930} is displayed as \texttt{0930}.

\section{Use the Module}

In a second file or Google Colab notebook, begin with

\begin{verbatim}
import course_tools
courses = []
\end{verbatim}

Add the courses you are taking this semester, and the courses of a classmate.  Use the format defined previously.

\section{Sort the Schedule}

Sort the courses by time:

\begin{verbatim}
course_tools.sort_by_time(courses)
course_tools.display_courses(courses)
\end{verbatim}

Predict which course will appear first and last.

\section{Challenge}

Extend the room-conflict function from the warm-up. Two courses should now be reported as a conflict only when they

\begin{itemize}
\item use the same building and room,
\item meet at the same time, and
\item share at least one meeting day.
\end{itemize}

For example, two courses in SLC 202 at 1400 should \textbf{not} conflict if one meets MWF and the other meets TR.

\end{document}